\documentclass[../../../../main.tex]{subfiles}

\begin{document}

\subsection*{第1問}

[解] $C = \cos t$, $S = \sin t$ とする。($0 \le t \le \pi$)
\[
\begin{cases}
x = 2\left(\frac{\sqrt{3}}{2} C + \frac{1}{2} S\right) = \sqrt{3} C + S \\
y = \frac{1}{2} C - \frac{\sqrt{3}}{2} S
\end{cases}
\]
だから, 題意の距離 $L$ として
\[
\begin{aligned}
L^2 &= x^2 + y^2 = \left( 3C^2 + S^2 + 2\sqrt{3}SC \right) + \left( \frac{1}{4}C^2 + \frac{3}{4}S^2 - \frac{\sqrt{3}}{2}SC \right) \\
&= \frac{13}{4} C^2 + \frac{7}{4} S^2 + \frac{3}{2}\sqrt{3} SC \\
&= \frac{7}{4} + \frac{3}{2} C^2 + \frac{3}{4}\sqrt{3} \sin 2t \\
&= \frac{5}{2} + \frac{3}{4} \{ \cos 2t + \sqrt{3} \sin 2t \} \\
&= \frac{5}{2} + \frac{3}{2} \sin(2t + \pi/6)
\end{aligned}
\]
より,
\[
\begin{cases}
t = \frac{2}{3}\pi \text{ で } \min L = 1 \\
t = \frac{\pi}{6} \text{ で } \max L = 2
\end{cases}
\]

\subsection*{第2問}

[解]
\[
\begin{cases}
a_{n+1} = \frac{1}{2}(a_n^2 + 1) \\
a_1 = \frac{1}{2}x
\end{cases}
\quad \cdots \text{①}
\]

(1) $a_{n+1} - a_n = \frac{1}{2}(a_n - 1)^2$ に注意する。漸化式から, 帰納的に $a_n > 0$ である。したがって, $a_n = 1$ なる $n \in \mathbb{N}$ があったと仮定すると, ①から $a_{n-1} = 1$ となり, 以下帰納的に $a_0 = 1 \iff x = 2$ となる。これは $x < 2$ に反し矛盾。よって $a_n \neq 1$ だから, 任意の $n$ に対し $a_{n+1} - a_n = \frac{1}{2}(a_n - 1)^2 > 0$ となるから, $a_1 < a_2 < \dots < a_n < \dots$ となる。 \quad [終]

(2) 題意を帰納法で示す。$n=1$ の時 $a_1 < \frac{1}{2} \cdot 2 = 1$ から成立。($\because x < 2$)
以下 $n=k \in \mathbb{N}$ での成立を仮定すると, ①から
\[
a_{k+1} = \frac{1}{2}(a_k^2 + 1) < \frac{1}{2}(1 + 1) = 1 \quad (\because 0 < a_k < 1)
\]
となって $n=k+1$ でも成立する。以上から任意の $n \in \mathbb{N}$ に対して $a_n < 1$ となる。 \quad [終]

次に後件部について考える。
\[
\varepsilon < 1 - \frac{1}{2}x, \quad x < 2 \quad \cdots \text{②}
\]
で, 題意および (1) から,
\[
a_n < 1 - \varepsilon < a_{n+1} \quad \cdots \text{③}
\]
(1) 及び ③ から,
\[
a_{k+1} - a_k = \frac{1}{2}(a_k - 1)^2 > \frac{1}{2} \left( (1 - \varepsilon) - 1 \right)^2 = \frac{1}{2} \varepsilon^2
\]
$k=1, 2, \dots, n$ として足して
\[
a_{n+1} - a_1 > \frac{1}{2} n \varepsilon^2
\]
$1 > a_{n+1}$, $a_1 = \frac{1}{2}x$ を代入して
\[
n \varepsilon^2 < 2 - x \quad \text{[終]}
\]
を得る。

\subsection*{第3問}

[解] $A = \int_0^1 \{f(x)\}^2 \, dx$ とする。又, 与式を $I$ とおく。
\[
\begin{aligned}
I &= A - 2 \int_0^1 f(x) (ax+b) \, dx + \int_0^1 (ax+b)^2 \, dx \\
&= A - 2 (3a + 2b) + \frac{1}{3} a^2 + ab + b^2 \\
&= b^2 + ab - 4b + \frac{1}{3} a^2 - 6a + A \\
&= \left( b + \frac{a-4}{2} \right)^2 + \frac{1}{12} a^2 - 4a + A - 4 \\
&= \left( b + \frac{a-4}{2} \right)^2 + \frac{1}{12} (a-24)^2 - 52 + A
\end{aligned}
\]
だから, $(a,b) = (24, -10)$ で $I$ は $\min$ となる。

\subsection*{第4問}

[解] $f_n(x) = 1 + \frac{x}{1!} + \dots + \frac{x^{n-1}}{(n-1)!}$ ($n \ge 2$) だから,
\[
f_n'(x) = f_{n-1}(x) \quad \text{[終]}
\]
となる。後半部について, $n=1, 2$ のとき,
\[
\begin{cases}
f_1(x) = 1 + x \\
f_2(x) = 1 + x + \frac{1}{2} x^2
\end{cases}
\]
だから, $n=1, 2$ では題意は成立する。そこで, $n=2k-1, 2k$ ($k \in \mathbb{N}$) での成立を仮定する。
\[
f_{2k+1}'(x) = f_{2k}(x) > 0 \quad (\because \text{仮定})
\]
から, $f_{2k+1}(x)$ は単調増加で, かつ $f_{2k+1}(x) \xrightarrow{x \to \pm \infty} \pm \infty$ だから, 中間値の定理より, $f_{2k+1}(x) = 0$ となる $x$ がただ1つある。これを $\alpha$ とする。($\alpha \neq 0$) $\quad \cdots \text{①}$
\[
f_{2k+2}'(x) = f_{2k+1}(x)
\]
から, 下表をうる。
\begin{center}
\begin{tabular}{c|c|c|c}
$x$ & $\dots$ & $\alpha$ & $\dots$ \\ \hline
$f'$ & $-$ & $0$ & $+$ \\ \hline
$f$ & $\searrow$ & & $\nearrow$
\end{tabular}
\end{center}
したがって,
\[
f_{2k+2}(x) \ge f_{2k+2}(\alpha) = \frac{\alpha^{2k+2}}{(2k+2)!} > 0 \quad (\because \alpha \neq 0, f_{2k+1}(\alpha) = 0)
\]
だから, $f_{2k+2}(x) = 0$ は実根をもたない $\quad \cdots \text{②}$

以上 ①, ② から, $n=2k+1, 2k+2$ でも成立。よって示された。 \quad [終]

\subsection*{第5問}

[解] $L(k) = |PQ_k|^2$ とおく。$\alpha_k = \frac{2k}{n}\pi$ とおくと,
\[
\begin{aligned}
L(k) &= (\cos \alpha_k - a)^2 + (\sin \alpha_k - b)^2 \\
&= a^2 + b^2 + 1 - 2 (a \cos \alpha_k + b \sin \alpha_k)
\end{aligned}
\]
だから,
\[
\begin{aligned}
S_n &= \frac{1}{n} \sum_{k=0}^{n-1} \left( a^2 + b^2 + 1 - 2(a \cos \alpha_k + b \sin \alpha_k) \right) \\
&\xrightarrow{n \to +\infty} a^2 + b^2 + 1 - 2 \int_0^1 (a \cos 2\pi x + b \sin 2\pi x) \, dx \\
&= a^2 + b^2 + 1
\end{aligned}
\]
より,
\[
\lambda_P = a^2 + b^2 + 1
\]
で, これは $a=b=0$ の時 $\min$ となる。

\subsection*{第6問}

[解1] まず $k \ge 2$ とする。$A \dots 3$ 人でじゃんけんをする。$B \dots 2$ 人でじゃんけんをする。とおくと,
$k$ 回目に1人の勝者が決まるのは以下の時:

\[
\begin{aligned}
&\text{1回目} \qquad \qquad a\text{回目} \qquad \qquad k\text{回目} \\
&\circ A \to A \to \dots \to A \to B \to \dots \to B \to 1\text{人} \quad \cdots \text{①} \\
&\circ A \to \dots \qquad \qquad \qquad \qquad \to A \to 1\text{人} \quad \cdots \text{②}
\end{aligned}
\]

$A$ を行う回数 $a$, $B$ を行う回数 $b$ とする。この時
\[
a+b = k \quad \cdots \text{③}
\]
である。

ここで, じゃんけんの推移についての確率を算出する。
\[
\begin{cases}
\circ 3\text{人} \to 2\text{人} \dots \underbrace{{}_3\mathrm{C}_1}_{\text{負ける人}} \cdot \underbrace{{}_3\mathrm{C}_1}_{\text{出す手}} \cdot \frac{1}{27} = \frac{1}{3} \\
\circ 3\text{人} \to 1\text{人} \dots {}_3\mathrm{C}_1 \cdot {}_3\mathrm{C}_1 \cdot \frac{1}{27} = \frac{1}{3} \\
\circ 3\text{人} \to 3\text{人} \dots 1 - \frac{1}{3} - \frac{1}{3} = \frac{1}{3} \quad (\text{あいこに注目}) \\
\circ 2\text{人} \to 1\text{人} \dots {}_2\mathrm{C}_1 \cdot {}_3\mathrm{C}_1 \cdot \frac{1}{9} = \frac{2}{3} \\
\circ 2\text{人} \to 2\text{人} \dots 1 - \frac{2}{3} = \frac{1}{3}
\end{cases}
\]

したがって, ①となる確率は,
\[
\left(\frac{1}{3}\right)^{a-1} \cdot \frac{1}{3} \cdot \left(\frac{1}{3}\right)^{b-1} \cdot \frac{2}{3} = 2 \cdot \left(\frac{1}{3}\right)^k \quad (\because \text{③})
\]
②となる確率は
\[
\left(\frac{1}{3}\right)^{k-1} \cdot \frac{1}{3} = \left(\frac{1}{3}\right)^k
\]
だから, $a = 1, 2, \dots, k-1$ として足して,
\[
\sum_{a=1}^{k-1} 2\left(\frac{1}{3}\right)^k + \left(\frac{1}{3}\right)^k = (2k-1)\left(\frac{1}{3}\right)^k \quad (k=1 \text{でも成立})
\]

\medskip

[解2] $n$ 回目に3人, 2人でジャンケンする確率を $a_n, b_n$ とする。又, もとめる確率 $c_k$ とすると,
\[
c_k = \frac{1}{3} a_k + \frac{2}{3} b_k \quad \cdots \text{①}
\]
である。又, $n, n+1$ 回目の推移から,
\[
\begin{cases}
a_{n+1} = \frac{1}{3} a_n \\
b_{n+1} = \frac{1}{3} a_n + \frac{1}{3} b_n
\end{cases}
\]
$a_1 = 1, b_1 = 0$ だから, これを解いて,
\[
a_n = \left(\frac{1}{3}\right)^{n-1}, \quad b_n = (n-1)\left(\frac{1}{3}\right)^{n-1}
\]
だから ① に代入して
\[
c_k = (2k-1)\left(\frac{1}{3}\right)^k
\]

\end{document}
